% !PRINTED 644
\weberpage{644}

As in § 175 one proves the theorems:

7. \gesperrt{Sum, difference, product of integral functionals in $\overline{\Omega}$ are again integral functionals.}

8. \gesperrt{If $\overline{\omega}$ is an integral functional in $\overline{\Omega}$ according to Definition 6. and at the same time rational, then it is an integral functional in $\overline{Z}$ (according to Definition 2.).}

9. \gesperrt{If $\overline{\omega}$ is an integral functional in $\overline{\Omega}$, then $N(\overline{\omega})$ is an integral functional in $\overline{Z}$, and the absolute of $N(\overline{\omega})$ is called the ``absolute norm of $\overline{\omega}$''.}

The absolute norm of $\overline{\omega}$ is therefore a function in $Z$.

If we denote the absolute norm by $N_a(\overline{\omega})$, then the theorem holds:
\weq{2}{N_a(\overline{\omega}\,\overline{\omega}') = N_a(\overline{\omega})\,N_a(\overline{\omega}').}

\begin{center}
§ 180. \gesperrt{Divisibility of functionals. Units.}
\end{center}

The theorems of Bd. II, § 155 are now to be repeated, with quite slight modifications, to which attention is drawn in what follows. To simplify the expression, integral functionals in $\overline{\Omega}$ shall here be denoted by the small Greek letters, integral functionals in $\overline{Z}$ by the small Latin letters. Then we have:

1. \gesperrt{Definition: If $\beta$ is different from zero, then $\alpha$ is called divisible by $\beta$ if $\alpha/\beta = \gamma$ is an integral functional in $\overline{\Omega}$.}

2. \gesperrt{If $\xi, \eta, \ldots$ are arbitrary integral functionals in $\overline{\Omega}$ and $\alpha, \beta, \ldots$ are divisible by $\delta$, then}
\[
\xi\alpha + \eta\beta + \cdots
\]
\gesperrt{is also divisible by $\delta$.}

3. \gesperrt{Definition: An integral functional $\varepsilon$ whose reciprocal $1/\varepsilon$ is integral, i.\,e. a divisor of the number 1, is called a unit in $\overline{\Omega}$. A unit is a divisor of every integral functional. Product and quotient of two units are again units.}

4. \gesperrt{Two integral functionals $\alpha$, $\beta$ which are mutually divisible by one another are called associated. Their quotient $\alpha/\beta = \varepsilon$ is a unit. Two functionals $\alpha$ and $\alpha\varepsilon$ are associated.}

5. \gesperrt{If $\alpha$ is divisible by $\beta$, then every functional associated with $\alpha$ is divisible by every functional associated with $\beta$.}
